% % % % % % % % % % % % % % % % % % % % % % % % % % % % % % % % % %
%  Devolutiva — Fase 1: Imersão
%  Disciplina: Design Thinking (DE_233) — ETE Cícero Dias
%  Profa. Samara Silva  ·  Recife, 1 de junho de 2026
% % % % % % % % % % % % % % % % % % % % % % % % % % % % % % % % % %

\documentclass[
  sigconf,
  language=portuguese
]{acmart}

% ── Desliga rodapé ACM ──────────────────────────────────────
\settopmatter{printacmref=false}
\setcopyright{none}
\acmYear{2026}

% ── Pacotes ─────────────────────────────────────────────────
\usepackage[portuguese]{babel}
\usepackage{microtype}
\usepackage{booktabs}
\usepackage{tabularx}
\usepackage{array}
\usepackage{colortbl}
\usepackage[dvipsnames,table]{xcolor}
\usepackage[inline]{enumitem}
\usepackage{graphicx}
\usepackage{adjustbox}

% ── Cores ───────────────────────────────────────────────────
\definecolor{cgood}{HTML}{1E6B3A}
\definecolor{cwarn}{HTML}{8B4500}
\definecolor{cred}{HTML}{C0392B}
\definecolor{cmuted}{HTML}{666666}
\definecolor{crow}{HTML}{F2EFEB}

% ── Colunas customizadas para tabela ────────────────────────
\newcolumntype{L}[1]{>{\raggedright\arraybackslash}p{#1}}
\newcolumntype{C}[1]{>{\centering\arraybackslash}p{#1}}

% ── Comando nota destacada ───────────────────────────────────
\newcommand{\notadestaque}[2]{%
  \noindent\textbf{Nota:}\enspace%
  {\large\bfseries\color{cred}#1\,/\,#2\,pts}%
}

% ── Comando label de status ──────────────────────────────────
\newcommand{\statusok}[1]{{\small\bfseries\color{cgood}#1}}
\newcommand{\statuswarn}[1]{{\small\bfseries\color{cwarn}#1}}

% ── Renomeações PT-BR ────────────────────────────────────────
\AtBeginDocument{%
  \renewcommand{\abstractname}{Resumo}
  \renewcommand{\figurename}{Figura}
  \renewcommand{\tablename}{Tabela}
}

% ────────────────────────────────────────────────────────────
\begin{document}

% ── CAPA ────────────────────────────────────────────────────
\thispagestyle{empty}
\IfFileExists{capa.png}{%
  \includegraphics[width=\paperwidth,height=\paperheight]{capa.png}%
}{%
  \begin{center}\small\color{gray}[capa.png n\~ao encontrada]\end{center}%
}
\newpage

% ── METADADOS ────────────────────────────────────────────────
\title{Devolutiva de Avaliação: Fase 1 --- Imersão}

\author{Profa.\ Samara Silva}
\authornote{Grupo avaliado: \textbf{Grupo 01 - Fome Zero} · Período: 22 de maio de 2026 · Avaliado em: 1 de junho de 2026.}
\email{samarasilvia.educa@gmail.com}
\affiliation{%
  \institution{ETE Cícero Dias}
  \department{Design Thinking --- DE\_233}
  \city{Recife}
  \state{Pernambuco}
  \country{Brasil}
}
\maketitle

\section{Avaliação Geral}

\notadestaque{3,27}{10,00}

A Fase~1 foi concluída. A seguir o resumo dos critérios e os comentários por critério.

\subsection{Resumo dos Critérios}

\rowcolors{2}{crow}{white}
\begin{tabular}{@{} L{2.8cm} C{2.0cm} C{1.1cm} C{1.1cm} @{}}
\toprule
\textbf{Critério} & \textbf{Status} & \textbf{Nota} & \textbf{Máx.} \\
\midrule
Pesquisa Desk & \statusok{Completo} & \textbf{0,50} & 0,50 \\
Matriz de Alinhamento & \statusok{Completo} & \textbf{0,50} & 0,50 \\
Matriz CSD & \statuswarn{Completo c/ ressalvas} & \textbf{0,42} & 0,50 \\
Pesquisa Primária & \statuswarn{Completo c/ ressalvas} & \textbf{0,85} & 1,00 \\
Relatório de Imersão & \statusok{Completo} & \textbf{1,00} & 1,00 \\
Persona com Empatia Real & \statuswarn{Não entregue} & \textbf{0,00} & 0,50 \\
Mapa de Empatia como Ferramenta & \statuswarn{Não entregue} & \textbf{0,00} & 0,50 \\
Coerência Problema → Solução & \statuswarn{Não entregue} & \textbf{0,00} & 0,75 \\
Jornada do Usuário & \statuswarn{Não entregue} & \textbf{0,00} & 0,25 \\
Geração de Ideias — Quantidade e Ousadia & \statuswarn{Não entregue} & \textbf{0,00} & 0,75 \\
Priorização com Critério & \statuswarn{Não entregue} & \textbf{0,00} & 0,75 \\
Inovação da Solução & \statuswarn{Não entregue} & \textbf{0,00} & 0,50 \\
Storytelling — A História do Usuário & \statuswarn{Não entregue} & \textbf{0,00} & 0,50 \\
Documentação do Processo & \statuswarn{Não entregue} & \textbf{0,00} & 0,75 \\
Colaboração e Distribuição do Trabalho & \statuswarn{Não entregue} & \textbf{0,00} & 0,75 \\
Participação no Pitch & \statuswarn{Não entregue} & \textbf{0,00} & 0,50 \\
\midrule
\textbf{Total} & & \textbf{\color{cred}3,27} & \textbf{10,00} \\
\bottomrule
\end{tabular}

\section{Comentários por Critério}

\subsection{Pesquisa Desk}

\noindent\statusok{0,50\,/\,0,50 pts · Completo}

\begin{itemize}[leftmargin=14pt,itemsep=1pt,topsep=3pt,parsep=0pt]
  \item[\textcolor{cgood}{$\checkmark$}] {\small\textcolor{cgood}{Dados secundários relevantes e atuais sobre o ODS escolhido}}
  \item[\textcolor{cgood}{$\checkmark$}] {\small\textcolor{cgood}{Fontes confiáveis: ONU, IBGE, artigos acadêmicos ou jornalísticos}}
  \item[\textcolor{cgood}{$\checkmark$}] {\small\textcolor{cgood}{Contextualiza o problema com base factual}}
  \item[\textcolor{cgood}{$\checkmark$}] {\small\textcolor{cgood}{Conexão clara entre os dados e o problema}}
\end{itemize}

O grupo está de parabéns pela qualidade da Desk Research! A escolha das fontes demonstra muita maturidade. Vocês buscaram referências com altíssima credibilidade (FAO, IBGE, EMBRAPA, MAPA) e dados extremamente atualizados, o que traz um embasamento fantástico para justificar a importância do projeto. Senti falta de dados diferentes como artigos ou algo de teor acadêmico, mas não tira o mérito do trabalho.

\textit{\textbf{Seção presente no relatório}}

A redação desta seção está excelente! O texto é extremamente claro, bem fundamentado e utiliza uma linguagem técnica e acadêmica perfeitamente adequada ao rigor exigido pelo projeto. A forma como vocês dividiram as frentes estratégicas (Análise Documental, Quantitativa e Levantamento Bibliográfico) demonstra muita organização e maturidade metodológica.

Outro ponto fortíssimo do trabalho foi a consolidação das inferências no final do texto. Extrair os dados da pesquisa teórica e estruturá-los de forma quantitativa (Dado 1 a Dado 6) para a construção de gráficos foi uma excelente decisão. Isso não apenas facilita a leitura e a visualização do problema, mas também comprova a capacidade analítica da equipe em transformar textos densos em informações acionáveis que justificam a criação do aplicativo.

Ótimo trabalho de síntese e argumentação! Parabéns, fiquei muito feliz com o trabalho de vocês.

\subsection{Matriz de Alinhamento}

\noindent\statusok{0,50\,/\,0,50 pts · Completo}

\begin{itemize}[leftmargin=14pt,itemsep=1pt,topsep=3pt,parsep=0pt]
  \item[\textcolor{cgood}{$\checkmark$}] {\small\textcolor{cgood}{Preenchida com perguntas reais do grupo sobre o problema}}
  \item[\textcolor{cgood}{$\checkmark$}] {\small\textcolor{cgood}{Reflete o que o grupo sabe e precisa descobrir}}
  \item[\textcolor{cgood}{$\checkmark$}] {\small\textcolor{cgood}{Serviu de base para planejar a pesquisa primária}}
\end{itemize}

Nada a declarar. Está tudo de acordo com o esperado!

\subsection{Matriz CSD}

\noindent\statuswarn{0,42\,/\,0,50 pts · Completo c/ ressalvas}

\begin{itemize}[leftmargin=14pt,itemsep=1pt,topsep=3pt,parsep=0pt]
  \item[\textcolor{cgood}{$\checkmark$}] {\small\textcolor{cgood}{Os três quadrantes preenchidos: Certezas, Suposições, Dúvidas}}
  \item[\textcolor{cgood}{$\checkmark$}] {\small\textcolor{cgood}{Conteúdo específico do projeto — não genérico}}
  \item[\textcolor{cgood}{$\checkmark$}] {\small\textcolor{cgood}{Revela pensamento crítico}}
  \item[\textcolor{cgood}{$\checkmark$}] {\small\textcolor{cgood}{Coerência entre a CSD e o que foi pesquisado}}
\end{itemize}

\subsection{Pesquisa Primária}

\noindent\statuswarn{0,85\,/\,1,00 pts · Completo c/ ressalvas}

\textit{\textbf{Seção presente no relatório}}

O texto que descreve os resultados da pesquisa está excelente. A redação é clara, objetiva e traz inferências muito ricas. A forma como vocês estruturaram os dados percentuais ao final (Dado Survey 1 ao 5) demonstra uma ótima capacidade de síntese. Vocês conseguiram cruzar as respostas para comprovar a viabilidade e a aceitação do aplicativo, o que atende perfeitamente ao objetivo desta etapa. Parabéns pela estruturação!

\textit{\textbf{InStrumento de coleta (Google forms)}}

Embora a estrutura básica do formulário esteja correta e os resultados tenham sido úteis, existem pontos importantes de melhoria para que a pesquisa alcance um rigor metodológico maior em projetos futuros:

\begin{enumerate}[leftmargin=14pt,itemsep=1pt,topsep=3pt]
  \item \textbackslash\{\}textbf\{Descrição e Termo de Consentimento:\} A introdução do formulário foi muito breve. Faltou especificar que vocês são estudantes do curso de Desenvolvimento de Sistemas da ETE Cícero Dias, detalhar um pouco mais o propósito do projeto e, muito importante em formulários, informar o tempo estimado de resposta (ex: "Esta pesquisa leva apenas 2 minutos").
  \item \textbackslash\{\}textbf\{Perfil do Usuário (Dados Demográficos):\} O formulário não coleta informações básicas de quem está respondendo. Mesmo sendo anônimo, é fundamental saber a faixa etária, a composição familiar (quantas pessoas moram na casa) e o contexto socioeconômico. Isso permite cruzar dados importantes (ex: "Será que famílias maiores desperdiçam mais?" ou "Qual a idade do público que mais aceita a tecnologia?").
  \item \textbackslash\{\}textbf\{Volume e Diversidade de Perguntas:\} Um questionário com apenas 5 perguntas limita muito o aprofundamento do tema. O ideal para uma pesquisa desse porte seria entre 10 e 12 perguntas, preferencialmente divididas por seções (Ex: Seção 1 - Perfil; Seção 2 - Hábitos de Consumo; Seção 3 - Relação com Tecnologia).
  \item \textbackslash\{\}textbf\{Tipos de Resposta:\} Houve uma limitação ao usar apenas perguntas de múltipla escolha (onde só se marca uma opção). Em pesquisas de campo, é enriquecedor incluir perguntas de "caixa de seleção" (onde o usuário pode marcar mais de um motivo) e, principalmente, perguntas abertas (texto curto ou longo) para captar opiniões espontâneas e insights qualitativos que a equipe não havia previsto.
\end{enumerate}
Apesar dessas oportunidades de melhoria no formulário, os dados coletados foram suficientes para embasar a justificativa do projeto nesta etapa. Excelente trabalho de análise!

\subsection{Relatório de Imersão}

\noindent\statusok{1,00\,/\,1,00 pts · Completo}

\begin{itemize}[leftmargin=14pt,itemsep=1pt,topsep=3pt,parsep=0pt]
  \item[\textcolor{cgood}{$\checkmark$}] {\small\textcolor{cgood}{Estrutura formal: introdução, resumo, metodologia, destaques, referências}}
  \item[\textcolor{cgood}{$\checkmark$}] {\small\textcolor{cgood}{Sintetiza achados da pesquisa primária e secundária}}
  \item[\textcolor{cgood}{$\checkmark$}] {\small\textcolor{cgood}{Há interpretação e conexão com o problema}}
  \item[\textcolor{cgood}{$\checkmark$}] {\small\textcolor{cgood}{Os achados alimentam a construção da persona}}
\end{itemize}

Toda a escrita, a articulação dos dados e a estrutura geral do relatório estão exemplares! O texto da Pesquisa Desk e a análise dos resultados do Survey foram redigidos com excelente maturidade acadêmica, clareza e uso correto de dados estatísticos. O trabalho demonstra um embasamento teórico e uma capacidade de síntese fantásticos. Parabéns pelo empenho na redação!

\textit{\textbf{Pontos de Revisão Crítica (Ajustes Obrigatórios)}}

Para garantir o rigor metodológico e evitar incoerências na entrega final, a equipe precisa corrigir dois problemas importantes na transição entre o planejamento e o relatório:

\begin{enumerate}[leftmargin=14pt,itemsep=1pt,topsep=3pt]
  \item \textbackslash\{\}textbf\{Ausência das Perguntas de Pesquisa no Relatório (Seção 2.2):\} Embora o título da seção seja \textbackslash\{\}textit\{"Matriz CSD e Perguntas de Pesquisa"\}, as perguntas em si não foram incluídas no corpo do relatório. Há uma quebra no texto que precisa ser corrigida para que o leitor entenda quais eixos guiaram a investigação.
\end{enumerate}
\textbf{Como corrigir:} > Na seção 2.2 do relatório, vocês devem incluir as verdadeiras \textbf{Perguntas de Pesquisa} (que são os objetivos estratégicos da \textit{equipe}, e não o questionário do público).

\subsection{Persona com Empatia Real}

\noindent\statuswarn{0,00\,/\,0,50 pts · Não entregue}

\begin{itemize}[leftmargin=14pt,itemsep=1pt,topsep=3pt,parsep=0pt]
  \item[$\square$] {\small\textcolor{cmuted}{A persona nasceu dos dados da imersão}}
  \item[$\square$] {\small\textcolor{cmuted}{Tem dor específica, não genérica}}
  \item[$\square$] {\small\textcolor{cmuted}{O grupo consegue explicar de onde veio cada característica}}
\end{itemize}

\subsection{Mapa de Empatia como Ferramenta}

\noindent\statuswarn{0,00\,/\,0,50 pts · Não entregue}

\begin{itemize}[leftmargin=14pt,itemsep=1pt,topsep=3pt,parsep=0pt]
  \item[$\square$] {\small\textcolor{cmuted}{Demonstra que o grupo tentou ver o mundo pela ótica do usuário}}
  \item[$\square$] {\small\textcolor{cmuted}{Não é só template preenchido — há reflexão real}}
  \item[$\square$] {\small\textcolor{cmuted}{O conteúdo alimentou a construção da persona}}
\end{itemize}

\subsection{Coerência Problema → Solução}

\noindent\statuswarn{0,00\,/\,0,75 pts · Não entregue}

\begin{itemize}[leftmargin=14pt,itemsep=1pt,topsep=3pt,parsep=0pt]
  \item[$\square$] {\small\textcolor{cmuted}{POV: [Persona] precisa de [necessidade] porque [insight]}}
  \item[$\square$] {\small\textcolor{cmuted}{O problema é específico e decorre da pesquisa}}
  \item[$\square$] {\small\textcolor{cmuted}{A solução responde diretamente à dor da persona}}
\end{itemize}

\subsection{Jornada do Usuário}

\noindent\statuswarn{0,00\,/\,0,25 pts · Não entregue}

\begin{itemize}[leftmargin=14pt,itemsep=1pt,topsep=3pt,parsep=0pt]
  \item[$\square$] {\small\textcolor{cmuted}{A jornada reflete o que foi descoberto na imersão}}
  \item[$\square$] {\small\textcolor{cmuted}{Identifica dores e oportunidades reais em cada etapa}}
\end{itemize}

\subsection{Geração de Ideias — Quantidade e Ousadia}

\noindent\statuswarn{0,00\,/\,0,75 pts · Não entregue}

\begin{itemize}[leftmargin=14pt,itemsep=1pt,topsep=3pt,parsep=0pt]
  \item[$\square$] {\small\textcolor{cmuted}{O brainstorming mostra que o grupo foi além do óbvio}}
  \item[$\square$] {\small\textcolor{cmuted}{Tem variedade de ideias, inclusive as ousadas}}
  \item[$\square$] {\small\textcolor{cmuted}{A ideação divergiu antes de convergir}}
\end{itemize}

\subsection{Priorização com Critério}

\noindent\statuswarn{0,00\,/\,0,75 pts · Não entregue}

\begin{itemize}[leftmargin=14pt,itemsep=1pt,topsep=3pt,parsep=0pt]
  \item[$\square$] {\small\textcolor{cmuted}{A matriz de impacto foi usada para tomar uma decisão real}}
  \item[$\square$] {\small\textcolor{cmuted}{Há justificativa para a escolha da solução final}}
  \item[$\square$] {\small\textcolor{cmuted}{O Canvas Lean demonstra que o grupo pensou sobre viabilidade}}
\end{itemize}

\subsection{Inovação da Solução}

\noindent\statuswarn{0,00\,/\,0,50 pts · Não entregue}

\begin{itemize}[leftmargin=14pt,itemsep=1pt,topsep=3pt,parsep=0pt]
  \item[$\square$] {\small\textcolor{cmuted}{O grupo identifica em qual tipo de inovação se encaixa sua solução}}
  \item[$\square$] {\small\textcolor{cmuted}{A justificativa é coerente com a solução apresentada}}
  \item[$\square$] {\small\textcolor{cmuted}{A solução não é a mais óbvia para aquele ODS}}
\end{itemize}

\subsection{Storytelling — A História do Usuário}

\noindent\statuswarn{0,00\,/\,0,50 pts · Não entregue}

\begin{itemize}[leftmargin=14pt,itemsep=1pt,topsep=3pt,parsep=0pt]
  \item[$\square$] {\small\textcolor{cmuted}{O pitch conta uma história — não apresenta slides sequencialmente}}
  \item[$\square$] {\small\textcolor{cmuted}{A narrativa segue: ODS → persona → dor real → solução → impacto}}
  \item[$\square$] {\small\textcolor{cmuted}{Quem assiste entende o problema antes de ver a solução}}
\end{itemize}

\subsection{Documentação do Processo}

\noindent\statuswarn{0,00\,/\,0,75 pts · Não entregue}

\begin{itemize}[leftmargin=14pt,itemsep=1pt,topsep=3pt,parsep=0pt]
  \item[$\square$] {\small\textcolor{cmuted}{Todas as fases estão documentadas e entregues no Drive}}
  \item[$\square$] {\small\textcolor{cmuted}{Dá pra rastrear a evolução do projeto}}
  \item[$\square$] {\small\textcolor{cmuted}{Há coerência entre as fases — cada etapa alimenta a próxima}}
\end{itemize}

\subsection{Colaboração e Distribuição do Trabalho}

\noindent\statuswarn{0,00\,/\,0,75 pts · Não entregue}

\begin{itemize}[leftmargin=14pt,itemsep=1pt,topsep=3pt,parsep=0pt]
  \item[$\square$] {\small\textcolor{cmuted}{O Trello mostra que o trabalho foi distribuído}}
  \item[$\square$] {\small\textcolor{cmuted}{O histórico de commits confirma participação de mais de um membro}}
  \item[$\square$] {\small\textcolor{cmuted}{Não tem um integrante fazendo tudo}}
\end{itemize}

\subsection{Participação no Pitch}

\noindent\statuswarn{0,00\,/\,0,50 pts · Não entregue}

\begin{itemize}[leftmargin=14pt,itemsep=1pt,topsep=3pt,parsep=0pt]
  \item[$\square$] {\small\textcolor{cmuted}{Todos os integrantes participam da apresentação}}
  \item[$\square$] {\small\textcolor{cmuted}{Cada membro demonstra conhecimento sobre o projeto como um todo}}
\end{itemize}

\section{Avaliação Geral}

\notadestaque{0,00}{10,00}

A Fase~1 foi concluída. A seguir o resumo dos critérios e os comentários por critério.

\subsection{Resumo dos Critérios}

\rowcolors{2}{crow}{white}
\begin{tabular}{@{} L{2.8cm} C{2.0cm} C{1.1cm} C{1.1cm} @{}}
\toprule
\textbf{Critério} & \textbf{Status} & \textbf{Nota} & \textbf{Máx.} \\
\midrule
Persona — Qualidade Visual e Estrutura & \statuswarn{Não entregue} & \textbf{0,00} & 0,75 \\
Jornada do Usuário — Qualidade Visual & \statuswarn{Não entregue} & \textbf{0,00} & 0,75 \\
Mapa de Empatia — Apresentação Visual & \statuswarn{Não entregue} & \textbf{0,00} & 0,50 \\
POV + Problema + Solução — Clareza & \statuswarn{Não entregue} & \textbf{0,00} & 0,50 \\
Documentos de Ideação — Qualidade Visual & \statuswarn{Não entregue} & \textbf{0,00} & 0,50 \\
Style Guide — Identidade Visual do Projeto & \statuswarn{Não entregue} & \textbf{0,00} & 0,50 \\
Gramática Visual — Princípios de Gestalt & \statuswarn{Não entregue} & \textbf{0,00} & 1,25 \\
Cor e Tipografia no Protótipo & \statuswarn{Não entregue} & \textbf{0,00} & 1,25 \\
Affordance, Feedback e Microinterações & \statuswarn{Não entregue} & \textbf{0,00} & 1,25 \\
Jornada e Fluxo de Tarefas no Protótipo & \statuswarn{Não entregue} & \textbf{0,00} & 1,25 \\
Progressão: Baixa → Média → Alta Fidelidade & \statuswarn{Não entregue} & \textbf{0,00} & 1,00 \\
Consistência Visual Geral & \statuswarn{Não entregue} & \textbf{0,00} & 0,50 \\
\midrule
\textbf{Total} & & \textbf{\color{cred}0,00} & \textbf{10,00} \\
\bottomrule
\end{tabular}

\section{Comentários por Critério}

\subsection{Persona — Qualidade Visual e Estrutura}

\noindent\statuswarn{0,00\,/\,0,75 pts · Não entregue}

\begin{itemize}[leftmargin=14pt,itemsep=1pt,topsep=3pt,parsep=0pt]
  \item[$\square$] {\small\textcolor{cmuted}{A persona é visual e bem diagramada — hierarquia clara e legível}}
  \item[$\square$] {\small\textcolor{cmuted}{Cor e tipografia coerentes}}
  \item[$\square$] {\small\textcolor{cmuted}{As informações estão organizadas com clareza}}
\end{itemize}

\subsection{Jornada do Usuário — Qualidade Visual}

\noindent\statuswarn{0,00\,/\,0,75 pts · Não entregue}

\begin{itemize}[leftmargin=14pt,itemsep=1pt,topsep=3pt,parsep=0pt]
  \item[$\square$] {\small\textcolor{cmuted}{Mapeia os pontos de contato: antes, durante e depois}}
  \item[$\square$] {\small\textcolor{cmuted}{Identifica dores e oportunidades de design}}
  \item[$\square$] {\small\textcolor{cmuted}{Apresentação visual clara com boa hierarquia}}
\end{itemize}

\subsection{Mapa de Empatia — Apresentação Visual}

\noindent\statuswarn{0,00\,/\,0,50 pts · Não entregue}

\begin{itemize}[leftmargin=14pt,itemsep=1pt,topsep=3pt,parsep=0pt]
  \item[$\square$] {\small\textcolor{cmuted}{Os quatro quadrantes preenchidos e organizados visualmente}}
  \item[$\square$] {\small\textcolor{cmuted}{A diagramação é clara — dá pra ler sem esforço}}
  \item[$\square$] {\small\textcolor{cmuted}{Há hierarquia visual entre os elementos}}
\end{itemize}

\subsection{POV + Problema + Solução — Clareza}

\noindent\statuswarn{0,00\,/\,0,50 pts · Não entregue}

\begin{itemize}[leftmargin=14pt,itemsep=1pt,topsep=3pt,parsep=0pt]
  \item[$\square$] {\small\textcolor{cmuted}{O POV está formulado de forma clara e direta}}
  \item[$\square$] {\small\textcolor{cmuted}{O problema é específico e bem descrito}}
  \item[$\square$] {\small\textcolor{cmuted}{A solução é coerente com o problema identificado}}
\end{itemize}

\subsection{Documentos de Ideação — Qualidade Visual}

\noindent\statuswarn{0,00\,/\,0,50 pts · Não entregue}

\begin{itemize}[leftmargin=14pt,itemsep=1pt,topsep=3pt,parsep=0pt]
  \item[$\square$] {\small\textcolor{cmuted}{Os documentos estão bem diagramados}}
  \item[$\square$] {\small\textcolor{cmuted}{Hierarquia visual clara: títulos, seções e conteúdo distinguíveis}}
\end{itemize}

\subsection{Style Guide — Identidade Visual do Projeto}

\noindent\statuswarn{0,00\,/\,0,50 pts · Não entregue}

\begin{itemize}[leftmargin=14pt,itemsep=1pt,topsep=3pt,parsep=0pt]
  \item[$\square$] {\small\textcolor{cmuted}{Paleta de cores com justificativa baseada em psicologia das cores}}
  \item[$\square$] {\small\textcolor{cmuted}{Tipografia escolhida com critério: hierarquia, legibilidade, caráter}}
  \item[$\square$] {\small\textcolor{cmuted}{Coerente com o ODS e com o tom da solução}}
\end{itemize}

\subsection{Gramática Visual — Princípios de Gestalt}

\noindent\statuswarn{0,00\,/\,1,25 pts · Não entregue}

\begin{itemize}[leftmargin=14pt,itemsep=1pt,topsep=3pt,parsep=0pt]
  \item[$\square$] {\small\textcolor{cmuted}{O layout aplica princípios de Gestalt de forma intencional}}
  \item[$\square$] {\small\textcolor{cmuted}{Proximidade: elementos relacionados estão agrupados}}
  \item[$\square$] {\small\textcolor{cmuted}{Similaridade: elementos com mesma função têm mesma aparência}}
  \item[$\square$] {\small\textcolor{cmuted}{Continuidade e/ou pregnância aplicados}}
\end{itemize}

\subsection{Cor e Tipografia no Protótipo}

\noindent\statuswarn{0,00\,/\,1,25 pts · Não entregue}

\begin{itemize}[leftmargin=14pt,itemsep=1pt,topsep=3pt,parsep=0pt]
  \item[$\square$] {\small\textcolor{cmuted}{A paleta do Style Guide foi aplicada com consistência}}
  \item[$\square$] {\small\textcolor{cmuted}{Contraste adequado entre texto e fundo}}
  \item[$\square$] {\small\textcolor{cmuted}{Hierarquia tipográfica clara: título, subtítulo, corpo, label}}
\end{itemize}

\subsection{Affordance, Feedback e Microinterações}

\noindent\statuswarn{0,00\,/\,1,25 pts · Não entregue}

\begin{itemize}[leftmargin=14pt,itemsep=1pt,topsep=3pt,parsep=0pt]
  \item[$\square$] {\small\textcolor{cmuted}{Botões e elementos interativos são visualmente reconhecíveis}}
  \item[$\square$] {\small\textcolor{cmuted}{O fluxo indica feedback claro após uma ação}}
  \item[$\square$] {\small\textcolor{cmuted}{Estados de hover/foco/ativo estão pensados}}
  \item[$\square$] {\small\textcolor{cmuted}{Microinterações ou transições foram consideradas}}
\end{itemize}

\subsection{Jornada e Fluxo de Tarefas no Protótipo}

\noindent\statuswarn{0,00\,/\,1,25 pts · Não entregue}

\begin{itemize}[leftmargin=14pt,itemsep=1pt,topsep=3pt,parsep=0pt]
  \item[$\square$] {\small\textcolor{cmuted}{O protótipo reflete a jornada do usuário mapeada}}
  \item[$\square$] {\small\textcolor{cmuted}{O fluxo de tarefas está claro — navegação lógica}}
  \item[$\square$] {\small\textcolor{cmuted}{O diagrama de telas mostra como as páginas se conectam}}
\end{itemize}

\subsection{Progressão: Baixa → Média → Alta Fidelidade}

\noindent\statuswarn{0,00\,/\,1,00 pts · Não entregue}

\begin{itemize}[leftmargin=14pt,itemsep=1pt,topsep=3pt,parsep=0pt]
  \item[$\square$] {\small\textcolor{cmuted}{Os três níveis foram entregues}}
  \item[$\square$] {\small\textcolor{cmuted}{O wireframe de baixa resolve estrutura e fluxo}}
  \item[$\square$] {\small\textcolor{cmuted}{O de média incorpora identidade visual}}
  \item[$\square$] {\small\textcolor{cmuted}{O de alta é o protótipo navegável final no Figma}}
\end{itemize}

\subsection{Consistência Visual Geral}

\noindent\statuswarn{0,00\,/\,0,50 pts · Não entregue}

\begin{itemize}[leftmargin=14pt,itemsep=1pt,topsep=3pt,parsep=0pt]
  \item[$\square$] {\small\textcolor{cmuted}{O projeto tem identidade visual coesa}}
  \item[$\square$] {\small\textcolor{cmuted}{Todos os documentos entregues têm coerência visual entre si}}
  \item[$\square$] {\small\textcolor{cmuted}{O protótipo é consistente com o Style Guide}}
\end{itemize}

\section{Avaliação Geral}

\notadestaque{0,00}{6,00}

A Fase~1 foi concluída. A seguir o resumo dos critérios e os comentários por critério.

\subsection{Resumo dos Critérios}

\rowcolors{2}{crow}{white}
\begin{tabular}{@{} L{2.8cm} C{2.0cm} C{1.1cm} C{1.1cm} @{}}
\toprule
\textbf{Critério} & \textbf{Status} & \textbf{Nota} & \textbf{Máx.} \\
\midrule
Estrutura e Completude & \statuswarn{Não entregue} & \textbf{0,00} & 0,50 \\
Qualidade da Escrita e Argumentação & \statuswarn{Não entregue} & \textbf{0,00} & 0,75 \\
Referências e Fontes & \statuswarn{Não entregue} & \textbf{0,00} & 0,25 \\
Protótipo Navegável e Completo & \statuswarn{Não entregue} & \textbf{0,00} & 0,50 \\
Fidelidade à Identidade Visual & \statuswarn{Não entregue} & \textbf{0,00} & 0,50 \\
Clareza e Fluidez da Apresentação & \statuswarn{Não entregue} & \textbf{0,00} & 0,50 \\
Slides — Qualidade Visual & \statuswarn{Não entregue} & \textbf{0,00} & 0,50 \\
Capacidade de Justificar Decisões & \statuswarn{Não entregue} & \textbf{0,00} & 0,50 \\
Organização e Uso Real do Trello & \statuswarn{Não entregue} & \textbf{0,00} & 0,75 \\
Sprint Reviews & \statuswarn{Não entregue} & \textbf{0,00} & 0,75 \\
Estrutura e Padrão de Nomenclatura & \statuswarn{Não entregue} & \textbf{0,00} & 0,50 \\
\midrule
\textbf{Total} & & \textbf{\color{cred}0,00} & \textbf{6,00} \\
\bottomrule
\end{tabular}

\section{Comentários por Critério}

\subsection{Estrutura e Completude}

\noindent\statuswarn{0,00\,/\,0,50 pts · Não entregue}

\begin{itemize}[leftmargin=14pt,itemsep=1pt,topsep=3pt,parsep=0pt]
  \item[$\square$] {\small\textcolor{cmuted}{O documento segue o template definido}}
  \item[$\square$] {\small\textcolor{cmuted}{Todas as seções estão preenchidas}}
  \item[$\square$] {\small\textcolor{cmuted}{O documento está em PDF com nomenclatura correta}}
\end{itemize}

\subsection{Qualidade da Escrita e Argumentação}

\noindent\statuswarn{0,00\,/\,0,75 pts · Não entregue}

\begin{itemize}[leftmargin=14pt,itemsep=1pt,topsep=3pt,parsep=0pt]
  \item[$\square$] {\small\textcolor{cmuted}{O texto é claro, objetivo e bem escrito}}
  \item[$\square$] {\small\textcolor{cmuted}{As decisões do projeto são justificadas}}
  \item[$\square$] {\small\textcolor{cmuted}{A argumentação é coerente com os documentos de DT e DCU}}
\end{itemize}

\subsection{Referências e Fontes}

\noindent\statuswarn{0,00\,/\,0,25 pts · Não entregue}

\begin{itemize}[leftmargin=14pt,itemsep=1pt,topsep=3pt,parsep=0pt]
  \item[$\square$] {\small\textcolor{cmuted}{Dados e afirmações têm fonte citada}}
  \item[$\square$] {\small\textcolor{cmuted}{Há pelo menos 3 fontes referenciadas}}
  \item[$\square$] {\small\textcolor{cmuted}{Referências formatadas de forma consistente}}
\end{itemize}

\subsection{Protótipo Navegável e Completo}

\noindent\statuswarn{0,00\,/\,0,50 pts · Não entregue}

\begin{itemize}[leftmargin=14pt,itemsep=1pt,topsep=3pt,parsep=0pt]
  \item[$\square$] {\small\textcolor{cmuted}{O protótipo está navegável no Figma — não são telas soltas}}
  \item[$\square$] {\small\textcolor{cmuted}{Cobre o fluxo principal com pelo menos 3 telas conectadas}}
  \item[$\square$] {\small\textcolor{cmuted}{Um usuário externo consegue navegar sem explicação}}
\end{itemize}

\subsection{Fidelidade à Identidade Visual}

\noindent\statuswarn{0,00\,/\,0,50 pts · Não entregue}

\begin{itemize}[leftmargin=14pt,itemsep=1pt,topsep=3pt,parsep=0pt]
  \item[$\square$] {\small\textcolor{cmuted}{O protótipo é consistente com o Style Guide}}
  \item[$\square$] {\small\textcolor{cmuted}{Cor, tipografia e componentes seguem o planejado}}
\end{itemize}

\subsection{Clareza e Fluidez da Apresentação}

\noindent\statuswarn{0,00\,/\,0,50 pts · Não entregue}

\begin{itemize}[leftmargin=14pt,itemsep=1pt,topsep=3pt,parsep=0pt]
  \item[$\square$] {\small\textcolor{cmuted}{O grupo apresenta com segurança — não lê do slide}}
  \item[$\square$] {\small\textcolor{cmuted}{Todos os integrantes participam}}
  \item[$\square$] {\small\textcolor{cmuted}{O tempo de 7 minutos é respeitado}}
\end{itemize}

\subsection{Slides — Qualidade Visual}

\noindent\statuswarn{0,00\,/\,0,50 pts · Não entregue}

\begin{itemize}[leftmargin=14pt,itemsep=1pt,topsep=3pt,parsep=0pt]
  \item[$\square$] {\small\textcolor{cmuted}{Os slides são visuais — não são paredes de texto}}
  \item[$\square$] {\small\textcolor{cmuted}{Hierarquia tipográfica clara}}
  \item[$\square$] {\small\textcolor{cmuted}{A identidade visual do projeto está presente nos slides}}
\end{itemize}

\subsection{Capacidade de Justificar Decisões}

\noindent\statuswarn{0,00\,/\,0,50 pts · Não entregue}

\begin{itemize}[leftmargin=14pt,itemsep=1pt,topsep=3pt,parsep=0pt]
  \item[$\square$] {\small\textcolor{cmuted}{O grupo responde perguntas sobre as escolhas com segurança}}
  \item[$\square$] {\small\textcolor{cmuted}{Sabe explicar por que fez o que fez}}
  \item[$\square$] {\small\textcolor{cmuted}{Quando não sabe, reconhece honestamente}}
\end{itemize}

\subsection{Organização e Uso Real do Trello}

\noindent\statuswarn{0,00\,/\,0,75 pts · Não entregue}

\begin{itemize}[leftmargin=14pt,itemsep=1pt,topsep=3pt,parsep=0pt]
  \item[$\square$] {\small\textcolor{cmuted}{O quadro tem as colunas corretas: Backlog, To Do, Doing, Done e Referências}}
  \item[$\square$] {\small\textcolor{cmuted}{Os cards têm título no formato verbo no imperativo}}
  \item[$\square$] {\small\textcolor{cmuted}{O histórico mostra uso ao longo do semestre}}
\end{itemize}

\subsection{Sprint Reviews}

\noindent\statuswarn{0,00\,/\,0,75 pts · Não entregue}

\begin{itemize}[leftmargin=14pt,itemsep=1pt,topsep=3pt,parsep=0pt]
  \item[$\square$] {\small\textcolor{cmuted}{Ao menos 3 Sprint Reviews foram preenchidos}}
  \item[$\square$] {\small\textcolor{cmuted}{Os reviews têm reflexão real: entregas, pendências, impedimentos e aprendizados}}
  \item[$\square$] {\small\textcolor{cmuted}{Há honestidade sobre o processo}}
\end{itemize}

\subsection{Estrutura e Padrão de Nomenclatura}

\noindent\statuswarn{0,00\,/\,0,50 pts · Não entregue}

\begin{itemize}[leftmargin=14pt,itemsep=1pt,topsep=3pt,parsep=0pt]
  \item[$\square$] {\small\textcolor{cmuted}{Todos os arquivos estão na pasta correta do Drive}}
  \item[$\square$] {\small\textcolor{cmuted}{Os arquivos seguem o padrão de nomenclatura}}
  \item[$\square$] {\small\textcolor{cmuted}{A indicação de turma (\_A ou \_B) está correta em todos}}
\end{itemize}

\section{Considerações Finais}

Os ajustes indicados são de natureza metodológica e devem ser
incorporados antes do avanço para a Fase~2 (Definição).
Continuem assim.

% ── Assinatura ───────────────────────────────────────────────

\noindent\rule{\linewidth}{0.4pt}

\begin{flushright}
\textbf{Profa.\ Samara Silva}\\
{\small\color{cmuted} Design Thinking --- DE\_233 · ETE Cícero Dias}\\
{\small\color{cmuted} Recife, 1 de junho de 2026}
\end{flushright}

% ── Sem referências bibliográficas ───────────────────────────

\end{document}
